\chapter{Inverse Dynamics}
\label{chap:5}
Inverse dynamics is the problem of finding the forces required to produce a given acceleration in a rigid-body system. Inverse dynamics calculations are used in motion control systems, trajectory design and optimization (e.g. for robots and animated figures), in mechanical design, and as a component in some forward dynamics algorithms. This last application is covered in Chapter 6. The inverse dynamics calculation can be summarized by the equation τ = ID(model, q, ˙q, ¨q) , where q, ˙q, ¨q and τ are vectors of generalized position, velocity, acceleration and force variables, respectively, and model is a system model of a particular rigid-body system. This chapter presents an algorithm to calculate the inverse dynamics of a kinematic tree. An algorithm for closed-loop systems appears in Chapter 8. The inverse dynamics of a kinematic tree is particularly easy to calculate, which is why we start with this problem. The algorithm is called the recursive NewtonEuler algorithm, and it is the simplest, most efficient algorithm for the job. It has a computational complexity of O(n), which means that the amount of calculation grows linearly with the number of bodies, or joint variables, in the tree. The algorithm is presented first in its basic form; then the implementation details are added; and then the spatial-vector version is compared with the original 3D vector version. This chapter also covers the topic of algorithm complexity, and explains why recursive algorithms are so efficient. 5.1 Algorithm Complexity We define the computational cost of an algorithm to be its operations count— the number of arithmetic operations performed in the course of executing the algorithm. If an algorithm involves real-number arithmetic, like a dynamics algorithm, then it is customary to count only the floating-point operations.

Thus, integer operations, like those required to calculate an array index or increment a loop variable, get ignored. The computational complexity of an algorithm is a measure of how the cost scales with the size of the problem. We use ‘big O’ notation to specify this quantity. If an algorithm is described as having O(np) complexity, where n is the size of the problem, then the cost grows in proportion to np as n →∞. To be more precise, an algorithm is said to have O(f(n)) complexity if there exist positive constants Cmin, Cmax and nmin such that the computational cost is at least Cminf(n), but no more than Cmaxf(n), for all n > nmin. The expression O(f(n)) can be interpreted as the set of all functions that grow in proportion to f(n) as n →∞. Thus, a statement like “Algorithm X is O(f(n))” means that the cost of X, expressed as a function of n, is an element of the set O(f(n)). An expression like O(f(n)) = O(g(n)) is then a set equality, and implies g(n) ∈O(f(n)) and vice versa. When quoting complexities, we select the simplest function in the set to serve as the argument to O. For example, if the exact cost of an algorithm is given by the polynomial a0 + a1n + · · · + apnp, then the simplest function that grows in proportion to this polynomial is np, and so we quote the complexity as O(np). In general, only the fastest-growing component of f matters. If f converges to a constant then we say that f is O(1). An algorithm’s complexity figure depends on how we measure the size of the problem. For example, the complexity of matrix inversion is said to be O(n3) because the cost of inverting an n × n matrix grows in proportion to n3 as n →∞. This figure relies on n being the dimension of the matrix. If n were instead the number of elements in the matrix, then the complexity figure would be O(n1.5). More generally, if m and n are two different measures of problem size, and m = g(n), then O(f(m)) = O(f(g(n))). For a dynamics algorithm, the size of the problem is the size of the rigidbody system. If the system is a kinematic tree, then there are two reasonable measures of system size: the number of bodies, NB, or the number of degrees of freedom, n. If every joint in the tree has at least one degree of freedom, then these two numbers satisfy NB ≤n ≤6NB, which implies that O(N p B) = O(np) . (5.1) There is therefore no need to distinguish between NB and n when quoting complexities. 5.2 Recurrence Relations Modern dynamics algorithms are recursive, and they are orders of magnitude faster than their non-recursive predecessors. They are called recursive because they make use of recurrence relations to calculate sequences of intermediate results. It is the use of recurrence relations that accounts for their efficiency; so let us take a moment to examine how these recurrence relations work.

A recurrence relation is a formula that defines the next element in a sequence in terms of previous elements. This formula, together with a set of initial values, provides a recursive definition of the sequence. For example, the sequence of Fibonacci numbers, 0, 1, 1, 2, 3, 5, 8, 13, . . . , is defined recursively by the recurrence relation Fi = Fi−1 + Fi−2 and the initial values F0 = 0 and F1 = 1. In a rigid-body system, if the bodies are numbered according to the rules in Section 4.1.2, then the sequence of body velocities, v0, v1, v2, . . . , is defined recursively by the recurrence relation vi = vλ(i) +Si ˙qi and a given initial value for v0. Let us examine the task of calculating body velocities, both with and without the use of recurrence relations. To keep things simple, let us assume that the system is an unbranched kinematic tree with 1-DoF joints. This means that λ(i) = i −1, and that the velocity across joint i is the product of the joint axis vector, si, with the scalar joint velocity variable ˙qi. The recurrence relation for body velocities is then vi = vi−1 + si ˙qi , (5.2) with the initial condition v0 = 0; and the non-recursive formula for body velocities is vi = i X j=1 sj ˙qj . (5.3) Let m and a denote the cost of a scalar-vector multiplication and a vector addition, respectively. The cost of one application of Eq. 5.2 is then m + a, and the cost of Eq. 5.3 is i m + (i −1)a. The cost of calculating the first n body velocities via Eq. 5.2 is therefore n(m + a), and the cost via Eq. 5.3 is 1 2(n(n + 1)m + n(n −1)a). Thus, the use of Eq. 5.2 results in an efficient O(n) calculation, while Eq. 5.3 leads to an inefficient O(n2) calculation. The source of the inefficiency is not hard to see. If we perform the task using Eq. 5.3, then we end up calculating v1 = s1 ˙q1 v2 = s1 ˙q1 + s2 ˙q2 ... vn = s1 ˙q1 + s2 ˙q2 + · · · + sn ˙qn . The product s1 ˙q1 is calculated n times over, when it only needed to be calculated once. Likewise, the product s2 ˙q2 is calculated n −1 times, the sum s1 ˙q1 +s2 ˙q2 is calculated n−1 times, and so on. Recurrence relations therefore offer a systematic way to avoid unnecessary repetition in calculating sequences of related quantities, and this is the secret of their efficiency. The computational advantage of recurrence relations can be even greater when the calculation is more complicated. For example, the recurrence relation for body accelerations in the unbranched chain is ai = ai−1 + si ¨qi + vi × si ˙qi , (5.4)

which is just the time-derivative of Eq. 5.2, and the non-recursive formula is ai = i X j=1 sj ¨qj + i X j=1 j−1 X k=1 (sk ˙qk) × (sj ˙qj) . (5.5) (In both cases, we have assumed that ˙si = vi × si.) Calculating the accelerations via Eq. 5.4 is an O(n) computation; but doing it via Eq. 5.5 is an O(n3) computation. Equations 5.3 and 5.5 express the body velocities and accelerations directly in terms of their elementary constituents—joint axis vectors and joint velocity and acceleration variables. Equations like this are said to be in closed form. The equation of motion of a rigid-body system, expressed in closed form, is τi = n X j=1 Hij ¨qi + n X j=1 n X k=1 Cijk ˙qj ˙qk + gi , (5.6) where Hij are the elements of the generalized inertia matrix, Cijk are the coefficients of the centrifugal and Coriolis forces, and gi are the gravity terms. Early, non-recursive algorithms calculated the dynamics via equations like this, and their computational complexities were typically O(n4). Recursive algorithms are far more efficient, and they have computational complexities as low as O(n). When the recursive Newton-Euler algorithm first appeared, it was found to be about a hundred times faster than its non-recursive predecessor, the Uicker/Kahn algorithm, for a 6-DoF robot arm (Hollerbach, 1980). Prior to the invention of efficient recursive algorithms, closed-form equations were the standard way of expressing and evaluating equations of motion. 5.3 The Recursive Newton-Euler Algorithm The inverse dynamics of a general kinematic tree can be calculated in three steps as follows: 1. Calculate the velocity and acceleration of each body in the tree. 2. Calculate the forces required to produce these accelerations. 3. Calculate the forces transmitted across the joints from the forces acting on the bodies. These are the steps of the recursive Newton-Euler algorithm. Let us now work out the equations for each step, assuming that the tree is modelled by the quantities λ(i), µ(i), Si and Ii, as explained in Chapter 4.

body i body λ(i) body j body k body l µ(i)=\{ j,k,l\} fi fi x fj joint i fl fk parent children Figure 5.1: Forces acting on body i Step 1: Let vi and ai be the velocity and acceleration of body i. The velocity of any body in a kinematic tree can be defined recursively as the sum of the velocity of its parent and the velocity across the joint connecting it to its parent. Thus, vi = vλ(i) + Si ˙qi . (v0 = 0) (5.7) (Initial values are shown in brackets.) The recurrence relation for accelerations is obtained by differentiating this equation, resulting in ai = aλ(i) + Si ¨qi + ˙Si ˙qi . (a0 = 0) (5.8) Successive iterations of these two formulae, with i ranging from 1 to NB, will calculate the velocity and acceleration of every body in the tree. Step 2: Let f B i be the net force acting on body i. This force is related to the acceleration of body i by the equation of motion f B i = Ii ai + vi ×∗Ii vi . (5.9) Step 3: Referring to Figure 5.1, let fi be the force transmitted from body λ(i) to body i across joint i, and let f x i be the external force (if any) acting on body i. External forces can model a variety of environmental influences: force fields, physical contacts, and so on. They are regarded as inputs to the algorithm; that is, their values are assumed to be known. The net force on body i is then f B i = fi + f x i − X j∈µ(i) fj , which can be rearranged to give a recurrence relation for the joint forces: fi = f B i −f x i + X j∈µ(i) fj . (5.10)

Successive iterations of this formula, with i ranging from NB down to 1, will calculate the spatial force across every joint in the tree. No initial values are required for this formula, as µ(i) is the empty set for any body with no children. Having computed the spatial force across each joint, it remains only to calculate the generalized forces at the joints, which are given by τi = ST i fi (5.11) (cf. Eq. 3.37). Equations 5.7 to 5.11 describe the recursive Newton Euler algorithm in its simplest form. External forces can be used to model a variety of environmental influences, including gravitational forces. However, a uniform gravitational field can be modelled more efficiently as a fictitious acceleration of the base. To do this, we simply replace the initial value for Eq. 5.8 with a0 = −ag, where ag is the gravitational acceleration vector. If this trick is used, then the calculated values of ai and f B i are no longer the true acceleration and net force for body i, but are offset by the gravitational acceleration and force vectors, respectively. (The gravitational force on body i is Ii ag.) Algorithm Features Equations 5.7 and 5.8 involve a propagation of data from the root to the leaves; Eq. 5.10 involves a propagation of data from the leaves to the root; and Eqs. 5.9 and 5.11 describe calculations that are local to each body. The recursive Newton-Euler algorithm is therefore a two-pass algorithm: it makes an outward pass through the tree, during which the velocities and accelerations are calculated, and then an inward pass, during which the joint forces are calculated. The calculation of body forces can be assigned to either the first or the second pass. The former is generally the better choice. This algorithm calculates more that just τi. Other quantities, like vi and fi, can sometimes be useful. For example, one of the tasks of a machine designer is to select joint bearings of sufficient strength to withstand the dynamic reaction forces during operation. These forces can be calculated from fi. By setting some inputs to zero, one can calculate specific components of the dynamics. For example, if the external force and joint velocity and acceleration terms are set to zero, but the fictitious base acceleration is retained, then the algorithm calculates only the gravity compensation terms—the terms that statically balance the force of gravity. Such terms are used in robot control systems. Alternatively, if we set gravity, acceleration and external force terms all to zero, but keep the joint velocities, then the algorithm calculates only the Coriolis and centrifugal terms. If efficiency is a concern, then a stripped-down version of the algorithm can be used to perform these specialized calculations. For example, if the objective is to calculate the gravity-compensation terms

only, then the algorithm simplifies to fi = −Ii ag + X j∈µ(i) fj (5.12) τi = ST i fi . (5.13) To calculate the actual gravity terms, replace −Ii ag with Ii ag. Equations in Body Coordinates Equations 5.7 to 5.11 express the algorithm without reference to any particular coordinate system. In practice, the algorithm works best if the calculations pertaining to body i are performed in body i coordinates (as defined in §4.2). This is the body-coordinates (or link-coordinates) version of the algorithm. Let us transform the equations of the recursive Newton-Euler algorithm into body coordinates. The quantities vi, ai, Si, Ii, fi and f B i are now all considered to be expressed in body i coordinates. Note that the system model already supplies Si and Ii in these coordinates. Expressed in body coordinates, Eqs. 5.7 and 5.8 become vi = iXλ(i) vλ(i) + Si ˙qi (v0 = 0) (5.14) and ai = iXλ(i) aλ(i) + Si ¨qi + ˚ Si ˙qi + vi × Si ˙qi . (a0 = −ag) (5.15) Compared with the original equations, the following changes have been made: a coordinate transform has been inserted to transform vλ(i) and aλ(i) from body λ(i) coordinates to body i coordinates; and the matrix ˙Si has been replaced with ˚ Si+vi×Si, where ˚ Si is the apparent derivative of Si in body coordinates. We have also taken the opportunity to replace the original initial acceleration value with a0 = −ag to simulate the effect of gravity. The joint calculation routine, jcalc, is able to supply the quantities Si, vJi = Si ˙qi and cJi = ˚ Si ˙qi, all in body i coordinates. We therefore modify Eqs. 5.14 and 5.15 as follows, in order to exploit the available data: vJi = Si ˙qi (5.16) cJi = ˚ Si ˙qi (5.17) vi = iXλ(i) vλ(i) + vJi (v0 = 0) (5.18) ai = iXλ(i) aλ(i) + Si ¨qi + cJi + vi × vJi . (a0 = −ag) (5.19) Equations 5.9 and 5.11 remain unchanged in the body-coordinates version, but Eq. 5.10 becomes fi = f B i −iX∗ 0 f x i + X j∈µ(i) iX∗ j fj . (5.20)

Basic Equations: v0 = 0 a0 = −ag vi = vλ(i) + Si ˙qi ai = aλ(i) + Si ¨qi + ˙Si ˙qi f B i = Ii ai + vi ×∗Ii vi fi = f B i −f x i + P j∈µ(i) fj τi = ST i fi Equations in Body Coordinates: v0 = 0 a0 = −ag vJi = Si ˙qi cJi = ˚ Si ˙qi vi = iXλ(i) vλ(i) + vJi ai = iXλ(i) aλ(i) + Si ¨qi + cJi + vi × vJi f B i = Ii ai + vi ×∗Ii vi fi = f B i −iX∗ 0 f x i + P j∈µ(i) iX∗ j fj τi = ST i fi Algorithm: v0 = 0 a0 = −ag for i = 1 to NB do [XJ, Si, vJ, cJ] = jcalc(jtype(i), qi, ˙qi) iXλ(i) = XJ XT(i) if λ(i)̸ = 0 then iX0 = iXλ(i) λ(i)X0 end vi = iXλ(i) vλ(i) + vJ ai = iXλ(i) aλ(i) + Si ¨qi + cJ + vi × vJ fi = Ii ai + vi ×∗Ii vi −iX∗ 0 f x i end for i = NB to 1 do τi = ST i fi if λ(i)̸ = 0 then fλ(i) = fλ(i) + λ(i)X∗ i fi end end Table 5.1: The recursive Newton-Euler equations and algorithm In this equation, we have assumed that the external forces emanate from outside the system, and have therefore assumed that they are expressed in absolute (i.e., body 0) coordinates. Equations 5.16 to 5.20, together with 5.9 and 5.11, are the equations of the body-coordinates version of the recursive Newton-Euler algorithm. Algorithm Details Table 5.1 shows the pseudocode for the body-coordinates version of the algorithm, together with the equations for both the basic version and the bodycoordinates version. Bear in mind that the symbols vi, ai, etc. in the bodycoordinates version are expressed in body coordinates, whereas the same symbols in the basic version represent either abstract vectors or coordinate vectors in an unspecified common coordinate system. The two-pass structure of the algorithm is evident from the two for loops in the pseudocode. The symbols XJ, vJ and cJ are local variables that take new values on each iteration of the first loop. If the tree contains joints for which the velocity variables are not the derivatives of the position variables, then

simply replace the symbols ˙qi and ¨qi with αi and ˙αi, respectively. Likewise, if the tree contains any joints that depend explicitly on time, then jcalc will automatically calculate vJ and cJ using Eqs. 3.32 and 3.41 instead of 5.16 and 5.17; and the only change in the pseudocode is to pass the current time as a fourth argument to jcalc. The only place where iX0 is used is to transform the external forces from base to body coordinates. If there are no external forces, then the if statement that calculates these transforms can be omitted. In implementing Eq. 5.20, we have transformed it from a calculation using µ(i) to a calculation using λ(i). A literal translation of Eq. 5.20 into pseudocode looks like this: for i = NB to 1 do fi = f B i −iX∗ 0 f x i for each j in µ(i) do fi = fi + iX∗ j fj end end However, the same calculation can be performed in a different order, but with the same end result, by the code fragment for i = 1 to NB do fi = f B i −iX∗ 0 f x i end for i = NB to 1 do if λ(i)̸ = 0 then fλ(i) = fλ(i) + λ(i)X∗ i fi end end In merging this code fragment with the rest of the algorithm, the body of the first loop becomes the last line in the body of the first loop in the algorithm. 5.4 The Original Version The recursive Newton-Euler algorithm was first described in a paper by Luh et al. (1980a), which was republished in Brady et al. (1982). This original algorithm was formulated using 3D vectors, and appears superficially to be quite different from the one described in Section 5.3. However, the two algorithms are actually almost the same, the only significant difference being that one uses spatial accelerations while the other uses classical accelerations (see §2.11). Table 5.2 shows a special case of the original algorithm in which the joints are all revolute. One of the drawbacks of using 3D vectors to express rigidbody dynamics is that the equations are sensitive to joint type. The algorithm

Original Equations: ωi = iEi−1 ωi−1 + z ˙qi vi = iEi−1 (vi−1 + ωi−1 × i−1ri) ˙ωi = iEi−1 ˙ωi−1 + z ¨qi + (iEi−1 ωi−1) × z ˙qi ˙vi = iEi−1( ˙vi−1 + ˙ωi−1 × i−1ri + ωi−1 × ωi−1 × i−1ri) Fi = mi ( ˙vi + ˙ωi × ci + ωi × ωi × ci) Ni = Icm i ˙ωi + ωi × Icm i ωi fi = Fi + iEi+1 fi+1 ni = Ni + ci × Fi + iEi+1ni+1 + iri+1 × iEi+1 fi+1 τi = zTni Symbols: ωi, ˙ωi angular velocity and acceleration of body i vi, ˙vi linear velocity and acceleration of the origin of Fi (body i coordinate frame) Fi, Ni net force and moment on body i, expressed at its centre of mass fi, ni force and moment exerted on body i through joint i mi, ci, Icm i mass, centre of mass, and rotational inertia about centre of mass for body i in Fi coordinates iEi−1 3 × 3 orthogonal rotation matrix transforming from Fi−1 to Fi coordinates i−1ri position of origin of Fi relative to Fi−1, expressed in Fi−1 coordinates ˙qi, ¨qi, τi joint i velocity, acceleration and force variables z direction of joint i axis in Fi coordinates Correspondence with version in Table 5.1: ˆvi = ωi vi  ˆai =  ˙ωi ˙vi −ωi × vi  ˆf B i = Ni + ci × Fi Fi  ˆfi = ni fi  iXi−1 =  iEi−1 0 −iEi−1 i−1ri× iEi−1  ˆIi = Icm i −mi ci× ci× mi ci× −mi ci× mi 1  Si =  z 0  z =

0 0 1

λ(i) = i −1 µ(i) = \{i + 1\} Table 5.2: The original recursive Newton-Euler equations

described in Luh et al. (1980a) was formulated for both revolute and prismatic joints, and it is full of equations of the form variable =  one expression if joint is revolute another expression if joint is prismatic. The choice of 3D vector symbols in this table is a compromise between the notation used in this book and that used in other books that list this algorithm. The equations in this table are written for an unbranched kinematic tree, in line with how this algorithm is described elsewhere, but it is a simple matter to modify them to work on a branched tree. Table 5.2 also shows the correspondence between the 3D vectors in the original algorithm and the spatial vectors in the algorithm shown in Table 5.1. There are only two places where the 3D vectors do not match the corresponding 3D component of a spatial vector: linear acceleration and net body moment. The latter is a trivial difference that is immediately rectified when Ni is used to calculate ni. It is caused by Ni being calculated at the centre of mass instead of at the body coordinate system’s origin. The former is a non-trivial difference, as it affects the quantities propagated from one body to the next.1 One interesting feature of the original algorithm is that the linear velocity is never used in any subsequent expression, and therefore need not be calculated. This is a special property that arises from using classical accelerations. If the velocity calculation is omitted, then this algorithm is slightly more efficient than the spatial version. However, the difference is only a few percent, and is easily dwarfed by other efficiency-related matters, such as the choice of programming language, the computer hardware, and so on. 5.5 Additional Notes The algorithm presented in Luh et al. (1980a) is the one that we usually regard as the recursive Newton-Euler algorithm. However, the idea of combining the Newton-Euler equations of motion with a recursive evaluation strategy predates this algorithm by a few years. In particular, one can see early forms of recursive Newton-Euler algorithms in Stepanenko and Vukobratovic (1976) and Orin et al. (1979). The algorithm has since been refined and optimized, and a few of these optimizations are covered in Chapter 10. Two of the fastest published versions can be found in He and Goldenberg (1989) and Balafoutis et al. (1988), the latter also appearing in Balafoutis and Patel (1991). At the time of 1According to Section 2.11, classical acceleration is the apparent derivative of spatial velocity in a coordinate system having a pure linear velocity equal to the velocity of the body-fixed point at the origin. In light of this, one can interpret the difference between the two algorithms as follows: the algorithm in Table 5.1 expresses the equations in stationary coordinate systems that coincide with the body coordinate frames at the current instant; whereas the algorithm in Table 5.2 expresses the equations in non-rotating coordinate systems that coincide with the body coordinate frames at the current instant, and that have the same linear velocities as the origins of the body coordinate frames.

its invention, Lagrangian dynamics was more popular than Newton-Euler; and this prompted Hollerbach to devise a recursive Lagrangian algorithm, and to compare the computational complexities of various recursive and nonrecursive algorithms (Hollerbach, 1980). This is the work that contrasted the O(n4) complexity of the nonrecursive approach, as typified by the Uicker/Kahn algorithm (Kahn and Roth, 1971), with the O(n) complexity of the recursive algorithms. The impetus for this kind of research within the robotics community at that time was to develop an inverse dynamics algorithm that would be fast enough for use in real-time control of robot motion. An early example of this can be found in Luh et al. (1980b), although the operational-space formulation eventually became more popular (Khatib, 1987). The relatively slow speed of computers at the time led some researchers to investigate parallel-computer implementations (Lathrop, 1985), while others investigated symbolic optimization techniques (Murray and Neuman, 1984; Neuman and Murray, 1987). One interesting new direction is the development of efficient algorithms to calculate the partial derivative of inverse dynamics with respect to a design parameter. An example of this can be found in Fang and Pollard (2003).
